\chapter{Topic: Casting and Run-Time Type Information}

\begin{aside}
Material from the remainder of the term, not covered in the original lecture
notes.
\end{aside}

\section{C-Style Casts}
In C:
\begin{lstlisting}[style=c]
Node n;
int *ip = (int *) &n;   // forces C to treat a Node* as an int*
\end{lstlisting}
C-style casts should be avoided in C++. They do several very different things
with the same syntax, they are hard to search for, and the compiler won't stop
you from doing something meaningless. If you must cast, use one of the four C++
casts, which say exactly what you mean.

\section{\code{static\_cast}}
For ``sensible'' casts with well-defined semantics, checked at compile time:
\begin{lstlisting}[style=cpp]
double d = 3.7;
void f(int x);
void f(double x);
f(static_cast<int>(d));    // calls f(int); d is truncated to 3
\end{lstlisting}
Downcasting a superclass pointer to a subclass pointer:
\begin{lstlisting}[style=cpp]
Book *b = new Text{...};
Text *t = static_cast<Text*>(b);
t->getTopic();
\end{lstlisting}
You are telling the compiler ``trust me, \code{b} really points to a Text''. If
you are wrong, the behaviour is undefined. \code{static\_cast} does no run-time
check.

\section{\code{reinterpret\_cast}}
For unsafe, implementation-dependent ``weird'' conversions: reinterpret the bits
of one type as another.
\begin{lstlisting}[style=cpp]
Student s;
Turtle *t = reinterpret_cast<Turtle*>(&s);   // treat a Student as a Turtle
\end{lstlisting}
Almost any use is undefined or non-portable. You will rarely, if ever, need it.

\section{\code{const\_cast}}
For converting between \code{const} and non-\code{const}; the only C++ cast
that can ``cast away'' const.
\begin{lstlisting}[style=cpp]
void g(int *p);               // suppose g doesn't actually modify *p,
                              // but its author forgot the const
void f(const int *p) {
    g(const_cast<int*>(p));   // OK as long as g really doesn't modify *p
}
\end{lstlisting}
Useful for working with libraries that are not const-correct. Modifying an
object that was originally declared const through such a pointer is undefined
behaviour.

\section{\code{dynamic\_cast}}
Is it safe to convert a \code{Book*} to a \code{Text*}? It depends on what the
pointer actually points to. \code{dynamic\_cast} checks at run time:
\begin{lstlisting}[style=cpp]
Book *pb = ...;
Text *pt = dynamic_cast<Text*>(pb);
if (pt) cout << pt->getTopic();   // pb really pointed to a Text
else cout << "Not a Text";        // dynamic_cast returned nullptr
\end{lstlisting}
\begin{itemize}
  \item If the cast works (\code{*pb} is a Text, or a subclass of Text),
        \code{pt} points to the object; otherwise \code{pt} is \code{nullptr}.
  \item \code{dynamic\_cast} only works on classes with at least one virtual
        method (it uses the vtable to find the run-time type --- see the next
        chapter). This information is called \textbf{RTTI} (run-time type
        information).
  \item With references: there's no null reference, so a failing
        \code{dynamic\_cast<Text\&>(b)} throws \code{std::bad\_cast}.
\end{itemize}

\subsection{Casting Smart Pointers}
\begin{lstlisting}[style=cpp]
shared_ptr<Book> sb = make_shared<Text>(...);
shared_ptr<Text> st = dynamic_pointer_cast<Text>(sb);   // shares ownership
// also: static_pointer_cast, const_pointer_cast
\end{lstlisting}

\section{Using RTTI: Polymorphic Assignment Revisited}
Dynamic casting lets us solve the mixed-assignment problem from Lecture 18 with a
virtual \code{operator=}:
\begin{lstlisting}[style=cpp]
Text &Text::operator=(const Book &other) {     // virtual, overrides Book's
    const Text &tother = dynamic_cast<const Text&>(other);  // throws bad_cast
    if (this == &tother) return *this;                      // if not a Text
    Book::operator=(other);
    topic = tother.topic;
    return *this;
}
\end{lstlisting}
Assigning a Comic to a Text through Book references now throws instead of
silently corrupting the Text.

\section{Is Dynamic Casting Good Style?}
\begin{lstlisting}[style=cpp]
void whatIsIt(Book *b) {
    if (dynamic_cast<Comic*>(b)) cout << "Comic";
    else if (dynamic_cast<Text*>(b)) cout << "Text";
    else if (b) cout << "Book";
    else cout << "Nothing";
}
\end{lstlisting}
This is tightly coupled to the Book hierarchy: every time a new subclass is
added, this function (and every function like it) must be found and changed. It
is exactly what virtual methods are for:
\begin{lstlisting}[style=cpp]
class Book {
  public:
    virtual std::string identify() const { return "Book"; }
};
// each subclass overrides identify()
\end{lstlisting}
When a hierarchy is \emph{closed} (the set of subclasses is fixed) but the set of
operations is open, Visitor (Lecture 20) is the virtual-method-based answer.
Code that switches on run-time type with \code{dynamic\_cast} is a design smell:
use virtual methods (or Visitor) whenever you can.

\section*{Summary}
\begin{center}
\begin{tabular}{@{}lp{0.62\textwidth}@{}}
\toprule
Cast & Use \\
\midrule
\code{static\_cast} & well-defined conversions; unchecked downcasts \\
\code{reinterpret\_cast} & reinterpret bits; unsafe, non-portable \\
\code{const\_cast} & add/remove \code{const} \\
\code{dynamic\_cast} & checked downcast using RTTI; \code{nullptr} or
                        \code{bad\_cast} on failure \\
\bottomrule
\end{tabular}
\end{center}
Avoid C-style casts; prefer virtual methods to type switches built from
\code{dynamic\_cast}.
